\documentclass[11pt,a4paper,UTF8,fontset=fandol]{ctexart}
\usepackage[margin=2.35cm]{geometry}
\usepackage{amsmath,amssymb,amsthm,mathtools,bm,booktabs,graphicx,hyperref,xcolor,enumitem,fancyhdr}
\usepackage{tikz}
\usetikzlibrary{arrows.meta,positioning,calc}
\tikzset{gene/.style={draw,rounded corners,minimum width=12mm,minimum height=8mm,fill=blue!6},act/.style={-{Stealth},thick,blue!65!black},rep/.style={-{Bar[width=6pt]},thick,red!70!black}}
\hypersetup{colorlinks=true,linkcolor=blue!50!black,urlcolor=blue!45!black,citecolor=blue!45!black}
\setlength{\headheight}{15pt}
\setlength{\parindent}{2em}\setlength{\parskip}{.25em}
\setcounter{tocdepth}{2}\numberwithin{equation}{section}
\newtheorem{proposition}{命题}[section]
\newtheorem{remark}{注}[section]
\newcommand{\dd}{\mathrm d}\newcommand{\R}{\mathbb R}\newcommand{\ssv}{\mathrm{ss}}
\pagestyle{fancy}\fancyhf{}\fancyhead[L]{最小基因调控回路：从反应到动力学}\fancyhead[R]{\thepage}
\title{\bfseries 最小基因调控回路的理论基础\\[4pt]\large 从化学反应到一、二、三组分系统及真实网络}
\author{可复现学习笔记（含源代码与数值图）}
\date{2026 年 10 月}
\begin{document}\maketitle
\begin{abstract}
本文采用明确的反应规则与动力学假设，从产生--降解、启动子结合和质量作用定律出发，研究少组分基因调控回路所允许的行为。主体依次涵盖单组分无反馈及自反馈、双组分单向和相互调控、三组分负反馈与前馈回路、出生--死亡随机过程；最后联系乳糖操纵子、组合调控、网络模体和模块嵌入。重点在于区分\emph{结构所允许的机制}、\emph{特定参数下实现的动力学}与\emph{具体生物学功能}，而非追求复杂网络的完整重构。
\end{abstract}
\newpage\tableofcontents\newpage
\section{范围、对象和约定}
基因调控网络的节点常取为基因表达产物（例如蛋白质）的有效浓度，并不意味着微观反应中只含这些分子。启动子状态、RNA 聚合酶、mRNA、配体、代谢物均可能在粗粒化中被消去。本文用 $x,y,z$ 表示无量纲浓度；正边表示偏导数 $\partial F_i/\partial x_j>0$，负边表示小于零，\emph{并不}指定具体结合机制。

标准连续模型为
\begin{equation}\dot x_i=F_i(\bm{x};\bm{u})-\gamma_i x_i,\qquad F_i\geq0,\quad\gamma_i>0.\end{equation}
其中 $F_i$ 为有效产生率，$\bm u$ 为外部输入。结论若只对无自反馈、瞬时调控或一阶降解成立，均明确标出；不将某一特殊参数结果当作结构本身的普遍性质。


\subsection{先检查解是否具有物理意义}
本文的确定性结论默认向量场局部 Lipschitz 连续，故给定初值有唯一局部解；涉及导数时再假设为 $C^1$。输入在讨论定态时固定。浓度空间为闭正象限，周期轨道的排除在其内部使用平面定理。
\begin{proposition}[正性与有界性]\label{prop:positive}
若 $F_i\geq0$，则非负初值产生的解保持非负。若进一步 $F_i\leq M_i$，则解存在于所有 $t\geq0$，并满足
\[
0\leq x_i(t)\leq x_i(0)e^{-\gamma_it}+\frac{M_i}{\gamma_i}(1-e^{-\gamma_it}).
\]
\end{proposition}
\begin{proof}
将 $\dot x_i+\gamma_ix_i=F_i(\bm x)$ 乘以积分因子并积分，得到
$x_i(t)=e^{-\gamma_it}x_i(0)+\int_0^t e^{-\gamma_i(t-s)}F_i(\bm x(s))\dd s$。
在解保持于正象限的区间内，右侧非负；边界 $x_i=0$ 上导数非负，正象限因而正向不变（局部唯一性排除了向外穿越）。用 $0\leq F_i\leq M_i$ 估计积分即得上界。有限时间内解留在有界闭集，局部解的延拓定理给出全局存在。
\end{proof}
\begin{figure}[htbp]\centering
\begin{tikzpicture}[node distance=22mm]
\node[gene] (x) {$X$}; \node[gene,right=of x] (y) {$Y$};
\draw[act] (x)--node[above]{激活：$+$}(y);
\node[gene,right=35mm of y] (u) {$X$}; \node[gene,right=of u] (v) {$Y$};
\draw[rep] (u)--node[above]{抑制：$-$}(v);
\end{tikzpicture}
\caption{调控边的记号。箭头表示产生率随调控变量增加，钝头表示产生率下降；边的符号不能单独指定结合机制。}\label{fig:signs}
\end{figure}


\section{从基本化学反应到速率方程}
\subsection{质量作用与产生--降解}
对于 $\sum_i\alpha_{ir}X_i\xrightarrow{k_r}\sum_i\beta_{ir}X_i$，质量作用律给出 $v_r(\bm x)=k_r\prod_i x_i^{\alpha_{ir}}$，以及
\begin{equation}\dot{\bm x}=S\bm v(\bm x),\qquad S_{ir}=\beta_{ir}-\alpha_{ir}.\end{equation}
其中 $S$ 是化学计量矩阵。对于 $\varnothing\xrightarrow{\alpha}X$, $X\xrightarrow{\gamma}\varnothing$，
\begin{equation}\dot x=\alpha-\gamma x,\qquad x(t)=\frac\alpha\gamma+\left(x(0)-\frac\alpha\gamma\right)e^{-\gamma t}.\end{equation}
唯一稳态 $x^*=\alpha/\gamma$ 全局稳定（对 $x\geq0$）。若增加 mRNA，
\begin{align}
\varnothing\xrightarrow{k_m}M,\quad M\xrightarrow{\delta_m}\varnothing,\quad M\xrightarrow{k_p}M+P,\quad P\xrightarrow{\delta_p}\varnothing,\\[-2pt]
\dot m=k_m-\delta_m m,\qquad\dot p=k_pm-\delta_p p.
\end{align}
仅当 mRNA 弛豫远快于蛋白质变化时，可近似令 $m\simeq k_m/\delta_m$；若研究振荡或短时间响应，该时间尺度消去可能失效。


\paragraph{两阶段表达的完整解与消去误差}
令 $m^*=k_m/\delta_m$、$p^*=k_pk_m/(\delta_p\delta_m)$。积分因子直接给出
\begin{align}
m(t)&=m^*+(m_0-m^*)e^{-\delta_mt},\\
p(t)&=p^*+(p_0-p^*)e^{-\delta_pt}
 +k_p(m_0-m^*)\frac{e^{-\delta_mt}-e^{-\delta_pt}}{\delta_p-\delta_m}.
\end{align}
当 $\delta_p=\delta_m=\delta$ 时，最后一项取极限为 $k_p(m_0-m^*)te^{-\delta t}$。各项均衰减，因此两阶段模型的唯一定态全局吸引。若转录率变为随时间变化的 $r(t)$，则
$m(t)=m_0e^{-\delta_mt}+\int_0^t e^{-\delta_m(t-s)}r(s)\dd s$。
若 $r$ 至少二阶可微且在所考察时间窗内缓变，在初始层 $t\gg1/\delta_m$ 之后，卷积积分的渐近展开为
$m(t)\simeq r(t)/\delta_m-\dot r(t)/\delta_m^2$。故仅有 $\delta_m\gg\delta_p$ 还不够，输入与其他调控变量也必须在 mRNA 时间尺度上缓慢变化。

\paragraph{无量纲化}
若物理浓度 $P$ 满足 $\dot P=A+B P^n/(K_P^n+P^n)-\gamma P$，取 $x=P/K_P$、$\tau=\gamma t$，则
$\dd x/\dd\tau=a+b x^n/(1+x^n)-x$，其中 $a=A/(\gamma K_P)$、$b=B/(\gamma K_P)$。后文的阈值应在对应单位下解读。


\subsection{启动子结合与调控函数}
考虑 $D+X\rightleftharpoons DX$，定义解离常数 $K_d=[D][X]/[DX]$。位点结合概率为 $p=x/(K_d+x)$。若未结合与已结合的转录速率分别为 $r_0,r_1$，则
\begin{equation}F(x)=r_0+(r_1-r_0)\frac{x}{K_d+x}.\label{eq:occupancy}\end{equation}
$r_1>r_0$ 表示激活，$r_1<r_0$ 表示抑制。这一式子需要结合状态相对表达动力学快速准平衡（或具有等效稳态占据）。Hill 表达式 $x^n/(K^n+x^n)$ 是有用的有效模型，但 $n$ 不应被直接解读为实际结合分子数。真实细菌启动子还可以出现多个操纵位点、协同结合和 DNA loop。


\paragraph{占据公式的推导}
令总启动子浓度 $D_T=[D]+[DX]$ 守恒。平衡条件
$k_{\rm on}[D]x=k_{\rm off}[DX]$ 给出 $K_d=k_{\rm off}/k_{\rm on}$，从而
$[DX]/D_T=x/(K_d+x)$；按两种状态的转录率加权即得式~\eqref{eq:occupancy}。
这里的 $x$ 是\emph{游离}调控因子浓度；若结合显著耗竭调控因子，必须同时解其守恒方程，不能直接把总浓度代入。
在 $x$ 暂时固定时，结合概率满足
$\dot p=k_{\rm on}x(1-p)-k_{\rm off}p$，其弛豫速率为 $k_{\rm on}x+k_{\rm off}$。该速率远大于表达变化速率时才支持准平衡近似。


\section{单组分系统：穷尽基本调控符号}

\begin{figure}[htbp]\centering
\begin{tikzpicture}
\node[gene] (none) at (0,0) {$X$};\node at (0,-.8) {无反馈};
\node[gene] (pos) at (4,0) {$X$};\node at (4,-.8) {自激活};
\node[gene] (neg) at (8,0) {$X$};\node at (8,-.8) {自抑制};
\draw[act] (pos) to[loop above,min distance=12mm] (pos);
\draw[rep] (neg) to[loop above,min distance=12mm] (neg);
\end{tikzpicture}
\caption{单节点的三种符号拓扑。所有模型仍包含产生与降解；自边表示自身浓度对产生率的作用。}
\end{figure}
\subsection{无反馈与单调自抑制}
考虑 $\dot x=F(x)-\gamma x$。若 $F(x)$ 连续、非负、严格递减且有界，则 $F(x)-\gamma x$ 严格递减；在合理边界条件下只有一个定态，且 $F'(x^*)-\gamma<0$。自抑制可以改变响应时间和随机涨落，却不会在该标量模型下产生多稳态。


\begin{proposition}[单调自抑制的全局收敛]
设 $F:[0,\infty)\to[0,\infty)$ 为 $C^1$、有界且非增，$\gamma>0$。则存在唯一平衡点 $x^*$，所有非负初值的解趋于它，且
$|x(t)-x^*|\leq e^{-\gamma t}|x(0)-x^*|$。
\end{proposition}
\begin{proof}
$f(x)=F(x)-\gamma x$ 严格递减，$f(0)\geq0$，且 $f(x)\to-\infty$，由介值定理有唯一零点（允许 $x^*=0$）。令 $e=x-x^*$，则
\[
\frac12\frac{\dd}{\dd t}e^2=e\{F(x)-F(x^*)\}-\gamma e^2\leq-\gamma e^2.
\]
积分得到所述估计；全局存在由命题~\ref{prop:positive}保证。若 $F'(x^*)<0$，小扰动的局部弛豫时间为 $1/(\gamma-F'(x^*))$，比相同降解率的无反馈模型更短。这是局部线性结论，不等于任意输入下的全程响应时间。
\end{proof}


\subsection{自激活与双稳态}
取
\begin{equation}\dot x=a+b\frac{x^n}{K^n+x^n}-\gamma x.\label{eq:auto}\end{equation}
定态是产生曲线与线性降解直线的交点。局部稳定条件为
\begin{equation}F'(x^*)-\gamma<0.\end{equation}
图~\ref{fig:auto} 的参数 $a=0.05,b=3.5,K=\gamma=1,n=4$ 产生三个定态，外侧稳定，中间不稳定。关键不是正反馈必然双稳态，而是正反馈与足够的非线性产生多个交点。临界折叠点满足 $F(x^*)=\gamma x^*$ 和 $F'(x^*)=\gamma$。
\begin{figure}[htbp]\centering\includegraphics[width=.74\linewidth]{figures/self_activation.pdf}\caption{自激活模型的产生率与降解率，实心点为稳定定态，空心点为不稳定定态；标注给出定态浓度。}\label{fig:auto}\end{figure}

\paragraph{如何检验双稳态，而非只观察交点}
对于光滑标量方程，若 $f'(x^*)<0$，则足够靠近定态时 $f$ 在左侧为正、右侧为负，轨道向定态运动；若 $f'(x^*)>0$ 则方向相反。导数为零时必须看高阶项，不能套用线性判据。对式~\eqref{eq:auto}，
\[
F'(x)=\frac{bnK^nx^{n-1}}{(K^n+x^n)^2}.
\]
当 $n>1$ 时该导数的唯一最大值出现在 $x_m=K((n-1)/(n+1))^{1/n}$：令 $h=x/K$，对 $h^{n-1}/(1+h^n)^2$ 求导，驻点条件为 $(n-1)-(n+1)h^n=0$。因此
\[
\max F'=\frac bK\frac{(n+1)^2}{4n}
\left(\frac{n-1}{n+1}\right)^{(n-1)/n}.
\]
若该值小于 $\gamma$，$f$ 严格递减，定态唯一；超过 $\gamma$ 只是出现多交点的必要条件，仍需选择合适的 $a$。
令 $q=(x^*/K)^n$，联立折叠条件可参数化为
\[
b_{\rm fold}=\frac{\gamma K}{n}q^{(1-n)/n}(1+q)^2,\qquad
a_{\rm fold}=\frac{\gamma K}{n}q^{1/n}(n-1-q).
\]
这是由 $F'=\gamma$ 先解出 $b$，再代入 $a=\gamma x-bq/(1+q)$ 得到。对 $a>0$ 需 $0<q<n-1$；非退化折叠还需 $f''\ne0$，而导数最大点处的退化情形需另作分析。图示参数确有三个简单根，而不只是采样中观察到三个交点：取 $f(x)=0.05+3.5x^4/(1+x^4)-x$，有 $f(0)>0$、$f(0.1)<0$、$f(1)>0$、$f(4)<0$，所以三个互不相交的区间 $(0,0.1)$、$(0.1,1)$、$(1,4)$ 各含至少一个根。上面的单峰导数计算说明 $f'$ 最多有两个零点，由 Rolle 定理 $f$ 最多有三个不同零点，故这三根已穷尽。两个极值分别严格低于、严格高于零，因此三个根均不在极值点处，斜率依次为负、正、负。数值位置为 $0.050022,0.686700,3.527541$，斜率见复现脚本输出。中间不稳定根将低、高表达态的吸引域分开，短暂扰动跨过此阈值可造成持久状态切换。


\begin{proposition}光滑的一维自治 ODE $\dot x=f(x)$ 不存在非平凡周期解。\end{proposition}
\begin{proof}
假设存在周期为 $T>0$ 的非恒定解。它在紧区间 $[0,T]$ 取得最大值；利用周期延拓，最大值时刻 $t_m$ 为局部极大点，故 $\dot x(t_m)=f(x(t_m))=0$。从初值 $x(t_m)$ 出发的常值函数也是解。局部唯一性迫使原解与此常值解一致，再延拓至整个周期，矛盾。
\end{proof}
显式延迟、外部周期输入或额外动态变量不属于该命题的假设；一维模型可以有双稳态，却不能自行维持周期振荡。

\section{双组分系统：符号分类与稳定性}
\subsection{单向调控}
\begin{equation}\dot x=\alpha_x-\gamma_xx,\qquad\dot y=G(x)-\gamma_yy.\end{equation}
$x$ 的定态先独立确定，再决定 $y$。纯前向连接不能让 $y$ 反过来影响 $x$；它改变输出幅度、响应时间和滤波性质，而不是构成闭合反馈回路。


\paragraph{定态、稳定性和动态滤波的推导}
若 $G$ 连续非负，则 $x^*=\alpha_x/\gamma_x$、$y^*=G(x^*)/\gamma_y$ 唯一。雅可比为下三角矩阵，特征值为 $-\gamma_x,-\gamma_y$。更强的全局结论来自
\[
y(t)=e^{-\gamma_yt}y_0+\int_0^t e^{-\gamma_y(t-s)}G(x(s))\dd s.
\]
由于 $x(s)\to x^*$，积分输入趋于 $G(x^*)$，将输入误差积分分成有限的早期区间与一致小的晚期区间，可得 $y(t)\to y^*$。若 $G$ 可微，小扰动 $\delta x$ 到 $\delta y$ 的频率响应为
$\delta y(i\omega)/\delta x(i\omega)=G'(x^*)/(\gamma_y+i\omega)$：高频幅值被抑制，相位发生滞后。这里把 $\delta x$ 当作指定输入，若输入是上游产生率扰动还要乘以上游的响应因子。


\subsection{双向调控：四种符号组合}
最简单的无自反馈模型是
\begin{equation}\dot x=f(y)-\gamma_xx,\qquad\dot y=g(x)-\gamma_yy.\label{eq:two}\end{equation}
两个交叉导数的符号给出 $++$（互激活）、$--$（互抑制）、$+-$ 和 $-+$（负反馈）。在定态，
\begin{equation}J=\begin{pmatrix}-\gamma_x&f'(y^*)\\g'(x^*)&-\gamma_y\end{pmatrix},\quad \operatorname{tr}J=-\gamma_x-\gamma_y<0,\quad\det J=\gamma_x\gamma_y-f'g'.\end{equation}
正反馈：$f'g'>0$。只有当 $f'g'>\gamma_x\gamma_y$ 时该定态变成鞍点；是否双稳态必须进一步分析全局零流线与非线性边界。负反馈：$f'g'<0$，每个定态局部稳定，且可能具有复特征值，产生衰减振荡。


\paragraph{二阶稳定判据的完整计算}
令 $q=f'(y^*)g'(x^*)$，则
\[
\lambda_{\pm}=-\frac{\gamma_x+\gamma_y}{2}
\pm\sqrt{\frac{(\gamma_x-\gamma_y)^2}{4}+q}.
\]
其特征多项式为 $\lambda^2+(\gamma_x+\gamma_y)\lambda+\gamma_x\gamma_y-q$。当行列式为正时，若根为实数，和为负且积为正迫使两根均负；若为共轭复数，实部为迹的一半，亦负。行列式为负时实根异号，故为鞍点。零行列式出现零特征值，线性化不足以判断稳定性。
负反馈的衰减振荡条件进一步是
$q<-(\gamma_x-\gamma_y)^2/4$，所以负反馈符号本身甚至不保证\emph{衰减}振荡；正反馈 $q>0$ 则两根一定为实数。
\begin{figure}[htbp]\centering
\begin{tikzpicture}
\foreach \i/\title in {0/互激活,1/互抑制,2/负反馈}{
\node at (4.1*\i+1,1.2) {\title};
\node[gene] (x\i) at (4.1*\i,0) {$X$};
\node[gene] (y\i) at (4.1*\i+2,0) {$Y$};}
\draw[act] (x0) to[bend left=25] (y0);\draw[act] (y0) to[bend left=25] (x0);
\draw[rep] (x1) to[bend left=25] (y1);\draw[rep] (y1) to[bend left=25] (x1);
\draw[act] (x2) to[bend left=25] (y2);\draw[rep] (y2) to[bend left=25] (x2);
\end{tikzpicture}
\caption{两个交叉边形成的反馈。互激活与互抑制均为正回路；一激活一抑制为负回路。交换节点得到另一种带标签的负回路。}
\end{figure}


\begin{proposition}对于模型~\eqref{eq:two}，在单连通正象限区域内不存在周期轨道。\end{proposition}
\begin{proof}
若存在非恒定周期轨道，解的唯一性使其成为简单闭曲线 $C$。设其围成区域为 $\Omega$，单连通性保证 $\Omega$ 仍在向量场光滑的区域内。向量场沿轨道与 $C$ 相切，故法向通量为零。由平面散度定理，
\[
0=\oint_C\bm v\cdot\bm n\,\dd\ell
=\iint_\Omega\nabla\cdot\bm v\,\dd x\dd y
=-(\gamma_x+\gamma_y)\operatorname{Area}(\Omega)<0,
\]
矛盾。这里 $\bm v=(f(y)-\gamma_xx,g(x)-\gamma_yy)$。
\end{proof}
注意此证明\emph{不适用于}任意二维负反馈系统：自调控使对角导数变化，非线性降解也改变散度。无周期轨道并不能单独排除多定态或证明全局收敛。

\begin{proposition}[双节点负反馈的唯一定态]
若 $f,g$ 为非负有界的 $C^1$ 函数，且一个严格递增、另一个严格递减，则式~\eqref{eq:two}有唯一非负定态，且局部渐近稳定。
\end{proposition}
\begin{proof}
消去 $y$，得到 $x=T(x):=f(g(x)/\gamma_y)/\gamma_x$。$T$ 非增且有界，所以 $T(x)-x$ 严格递减，在零处非负，充分大时为负，故恰有一个根；$y=g(x)/\gamma_y$ 随之唯一。交叉斜率乘积 $q\leq0$，由上述判据其行列式严格为正、迹严格为负，给出局部稳定。
\end{proof}

\subsection{互相抑制与双稳态开关}
典型无量纲模型：
\begin{equation}\dot x=\frac{a}{1+y^n}-x,\qquad\dot y=\frac{a}{1+x^n}-y.\label{eq:toggle}\end{equation}
零流线 $x=a/(1+y^n)$ 与 $y=a/(1+x^n)$ 的交点就是定态。图~\ref{fig:toggle} 展示 $a=3,n=3$ 的示例：两个不对称表达态和中间的对称定态。可通过雅可比矩阵逐一判定稳定性。该模型展示状态选择的最小机制，但生物学中的细胞命运并不等同于某一对基因浓度。


\paragraph{对称定态失稳阈值}
令 $x=y=s$，则 $a=s(1+s^n)$；右侧在 $s\geq0$ 严格递增，故对称定态唯一。记交叉斜率的绝对值为
$h=ns^n/(1+s^n)$，雅可比为 $\begin{psmallmatrix}-1&-h\\-h&-1\end{psmallmatrix}$。
对称扰动 $(1,1)$ 对应特征值 $-1-h$，反对称扰动 $(1,-1)$ 对应 $-1+h$。当 $n>1$ 时，后者过零的阈值为
\[
s_c=(n-1)^{-1/n},\qquad a_c=\frac n{n-1}(n-1)^{-1/n}.
\]
当 $n\leq1$ 时有限 $s>0$ 下 $h<1$，对称态不会经这一线性机制失稳。阈值只判断对称态；不对称态的存在和稳定性仍须证明。

\begin{proposition}[图中参数确有两个稳定态]
式~\eqref{eq:toggle}在 $a=3,n=3$ 时恰有三个非负定态：一个对称鞍点和两个不对称的局部渐近稳定点。
\end{proposition}
\begin{proof}
定态满足 $x+xy^3=3$、$y+yx^3=3$，故 $x,y>0$。相减得
$(x-y)[1-xy(x+y)]=0$。对称解满足 $s+s^4=3$，严格单调性保证唯一；又 $s>1$，故 $h=3s^3/(1+s^3)>1$，是鞍点。
不对称时令 $S=x+y$、$P=xy$，有 $PS=1$。相加得
$S+P(S^2-2P)=6$，即 $S-S^{-2}=3$。左侧在 $S>0$ 严格递增，故恰有一个解，且 $S>3$。二次方程 $t^2-St+1/S=0$ 的判别式 $S^2-4/S>0$，两根均正；交换两根得恰好两个不对称解。此和、差条件亦反推出原两条定态方程，未引入额外解。
在这些点，两个交叉斜率均等于 $-3x^2y^2/a=-P^2$，特征值为 $-1\pm P^2$。因为 $P=1/S<1/3$，两根均负。三种分支已穷尽所有定态。
\end{proof}

\begin{figure}[htbp]\centering\includegraphics[width=.67\linewidth]{figures/toggle_nullclines.pdf}\caption{互抑开关的零流线与相流（$a=3,n=3$）。实心点稳定，空心点为鞍点；绿色轨迹分别趋于两种不对称表达态。}\label{fig:toggle}\end{figure}

\subsection{自调控使分类增多}
若允许 $x\to x$ 与 $y\to y$ 的正、负或缺失边，共有 $3^4=81$ 种带标签的二节点符号拓扑（四条可能有向边，每条三种取值）。这是\emph{拓扑计数}，不是 81 种不等价动力学机制；交换节点可合并部分结构，而固定符号也不唯一确定启动子逻辑、时延和参数。应先掌握无自调控的解析分类，再考察自调控如何改变迹、行列式与 Bendixson 条件。


\paragraph{拓扑计数的证明与节点交换}
四条边 $(X\to X,X\to Y,Y\to X,Y\to Y)$ 各有缺失、正、负三种独立取值，乘法原理给出 $3^4=81$。若只按交换 $X,Y$ 合并，交换后不变的图必须有两条自边相同、两条交叉边相同，故有 $3^2=9$ 个。其余 $81-9$ 个两两配对，得到 $9+(81-9)/2=45$ 个无标签符号拓扑。这个计数包含不连通图；若另要求连通或无自边，需要重新筛选。


\section{三组分：为什么增加一个环节有意义}
\subsection{负反馈环与振荡}
三个相互抑制环节可写为（本节取相同的一阶降解率 $\gamma$）
\begin{equation}\dot x_1=f_1(x_3)-\gamma x_1,\quad\dot x_2=f_2(x_1)-\gamma x_2,\quad\dot x_3=f_3(x_2)-\gamma x_3.\end{equation}
若平衡点处 $f'_1f'_2f'_3=-k^3<0$，特征方程为
\begin{equation}(\lambda+\gamma)^3=-k^3.\end{equation}
因此 $\lambda_1=-\gamma-k$ 与 $\lambda_{2,3}=-\gamma+k/2\pm i\sqrt3k/2$，在 $k=2\gamma$ 时一对特征值过虚轴。这给出线性失稳门槛；真正的稳定极限环还需要检查非线性、边界有界性及分岔条件。图~\ref{fig:rep} 使用 Hill 抑制 $f(s)=20/(1+s^3)$ 作数值说明，不能由图代替定理。


\paragraph{特征方程和阈值的推导}
写 $a_i=f_i'$，则
\[
\lambda I-J=\begin{pmatrix}\lambda+\gamma&0&-a_1\\-a_2&\lambda+\gamma&0\\0&-a_3&\lambda+\gamma\end{pmatrix},
\quad \det(\lambda I-J)=(\lambda+\gamma)^3-a_1a_2a_3.
\]
当乘积为 $-k^3$ 时，取 $-1$ 的三个立方根 $-1,\frac12\pm i\frac{\sqrt3}{2}$ 即得前述特征值。$k<2\gamma$ 时三根实部均负，$k>2\gamma$ 时共轭对实部为正，临界角频率为 $\sqrt3\gamma$。
对对称 Hill 环 $f(s)=\alpha/(1+s^n)$，唯一对称定态满足 $\alpha=\gamma s(1+s^n)$，因此
\[
\frac k\gamma=\frac{ns^n}{1+s^n}.
\]
事实上环中三个函数均相同且严格递减时，消元后的三次复合映射严格递减，固定点唯一，所以这个对称解也是整个环的唯一定态。
当 $n\leq2$ 时不可能达到有限参数的阈值；当 $n>2$ 时
\[
s_c=\left(\frac2{n-2}\right)^{1/n},\qquad
\alpha_c=\gamma\frac n{n-2}\left(\frac2{n-2}\right)^{1/n}.
\]
在 $n=3,\gamma=1$ 下 $\alpha_c=3\sqrt[3]2\approx3.780$，演示的 $\alpha=20$ 位于定态失稳一侧。
随 $\alpha$ 增加，$\dd s/\dd\alpha=1/[\gamma(1+(n+1)s^n)]>0$，且 $\dd k/\dd s>0$；故共轭对横截虚轴。然而要引用非退化 Hopf 分岔定理，仍需检验非零第一 Lyapunov 系数；其符号决定小振幅周期解的稳定性。本文不宣称已证明稳定极限环，所示持续振荡为给定初值与参数的数值结果。
\begin{figure}[htbp]\centering
\begin{tikzpicture}
\node[gene] (a) at (0,1.4) {$X_1$};\node[gene] (b) at (3,1.4) {$X_2$};\node[gene] (c) at (1.5,-.5) {$X_3$};
\draw[rep] (a)--(b);\draw[rep] (b)--(c);\draw[rep] (c)--(a);
\node[gene] (u) at (6,1.4) {$u$};\node[gene] (y) at (9,1.4) {$y$};\node[gene] (z) at (7.5,-.5) {$z$};
\draw[act] (u)--(y);\draw[act] (u)--node[left]{慢支路}(z);\draw[rep] (z)--(y);
\end{tikzpicture}
\caption{左：三次抑制构成负反馈环。右：不一致前馈回路，直接激活与间接抑制在输出处竞争；输入 $u$ 在这里是外部信号。}
\end{figure}

\begin{figure}[htbp]\centering\includegraphics[width=.78\linewidth]{figures/repressilator.pdf}\caption{三节点负反馈环的数值轨迹，生成脚本见 \texttt{code/make\_figures.py}。}\label{fig:rep}\end{figure}

\subsection{不一致前馈与适应}
输入 $u$ 可以直接激活输出 $y$，又通过慢变量 $z$ 抑制 $y$：
\begin{equation}\tau_y\dot y=H(u,z)-y,\qquad\tau_z\dot z=G(u)-z.\end{equation}
若 $z$ 比 $y$ 慢，阶跃输入引起 $y$ 先增加，随后 $z$ 增加抑制 $y$，形成脉冲。\emph{脉冲}与\emph{精确适应}不同：精确适应要求稳态输出 $y^*(u)$ 与持续输入的幅度无关，这对 $H$、$G$ 有额外条件；拓扑本身不保证这种性质。Ma 等（2009）对三节点生化适应网络进行系统枚举，识别出不一致前馈和带缓冲节点负反馈等机制，但其研究对象并不是所有可能的三基因转录网络。


\begin{proposition}[前馈模型的精确适应条件]
对固定输入区间 $I$，若 $H,G$ 连续且对应解收敛，则其稳态输出为 $y^*(u)=H(u,G(u))$。精确适应当且仅当该复合函数在 $I$ 上为常数。若函数为 $C^1$ 且 $I$ 为连通区间，该条件等价于
\[
H_u(u,G(u))+H_z(u,G(u))G'(u)=0\quad(u\in I).
\]
\end{proposition}
\begin{proof}
固定输入的第二个方程解为 $z(t)=G(u)+(z_0-G(u))e^{-t/\tau_z}$，故 $z\to G(u)$。第一个方程的积分因子公式与单向调控相同，连续输入 $H(u,z(t))\to H(u,G(u))$ 推出 $y$ 收敛于该值。输入无关恰是复合函数恒定；链式法则与导数为零的常函数判据给出后一等价式。
\end{proof}
\paragraph{一个可计算的适应例子}
取 $u>0,z>0$、$G(u)=u$、$H(u,z)=u/z$。每个恒定输入的定态为 $(y^*,z^*)=(1,u)$，故确有精确适应。若从输入 $u_0$ 的稳态出发，在 $t=0$ 切换至 $u_1>u_0$，则
\[
z(t)=u_1+(u_0-u_1)e^{-t/\tau_z},\qquad
\dot y(0^+)=\frac{u_1/u_0-1}{\tau_y}>0.
\]
$y$ 起初升高，最终回到 $1$，所以至少存在一个高于基线的峰；$\tau_z\gg\tau_y$ 使脉冲更易分辨。$u/z$ 是理想化的比值调控，只在 $u,z>0$ 的工作域内使用。若换成 $H=u/(1+z)$，保持同一拓扑与 $G$，则 $y^*=u/(1+u)$ 依赖输入，精确适应消失。这给出了拓扑不充分的明确反例。
\begin{figure}[htbp]\centering\includegraphics[width=.85\linewidth]{figures/feedforward_adaptation.pdf}
\caption{相同的不一致前馈结构可以精确适应，也可以只产生瞬态峰而改变最终基线。输入从 $1$ 跳至 $2$，$\tau_y=0.25,\tau_z=2$；右图比较两种调控函数的稳态输入输出曲线。}\label{fig:adapt}
\end{figure}


\section{随机化学反应：从平均值到分布}
出生--死亡过程
\begin{equation}\varnothing\xrightarrow{k}X,\qquad X\xrightarrow{\gamma}\varnothing\end{equation}
的主方程是
\begin{equation}\dot P_n=kP_{n-1}+\gamma(n+1)P_{n+1}-(k+\gamma n)P_n,\quad P_{-1}=0.\end{equation}
稳态满足相邻态详细平衡关系 $kP_{n-1}=\gamma nP_n$，故
\begin{equation}P_n=e^{-\mu}\frac{\mu^n}{n!},\qquad\mu=k/\gamma,\qquad\operatorname{Var}(n)=\langle n\rangle=\mu.\end{equation}


\paragraph{主方程为何具有这些项}
在短时间 $\dd t$ 内，$n-1\to n$ 的概率率为 $k$，$n+1\to n$ 的概率率为 $\gamma(n+1)$；从 $n$ 离开的总率为 $k+\gamma n$。只保留 $\dd t$ 一阶项，流入减流出即得主方程，$n=0$ 时没有向负状态的跃迁。
\begin{proposition}[泊松稳态、矩和相关时间]
$k>0,\gamma>0$ 的出生--死亡链的唯一归一化稳态为上述泊松分布。具有有限二阶矩时，
\begin{align}
\frac{\dd\langle n\rangle}{\dd t}&=k-\gamma\langle n\rangle,\\
\frac{\dd\operatorname{Var}(n)}{\dd t}&=k+\gamma\langle n\rangle-2\gamma\operatorname{Var}(n).
\end{align}
稳态自协方差为 $\operatorname{Cov}(n(t),n(0))=\mu e^{-\gamma t}$，$t\geq0$。
\end{proposition}
\begin{proof}
定义从 $n$ 到 $n+1$ 的概率流 $j_n=kP_n-\gamma(n+1)P_{n+1}$。稳态主方程为 $j_{n-1}-j_n=0$，而零边界满足 $j_{-1}=0$，故所有 $j_n=0$。递推得 $P_n=P_0\mu^n/n!$；归一化 $\sum_n\mu^n/n!=e^\mu$ 决定 $P_0=e^{-\mu}$，也证明唯一性。
对可积函数 $h(n)$，跃迁规则给出
\[
\frac{\dd\langle h(n)\rangle}{\dd t}
=\left\langle k[h(n+1)-h(n)]+\gamma n[h(n-1)-h(n)]\right\rangle.
\]
取 $h=n$ 得均值方程；取 $h=n^2$ 得
$\dd\langle n^2\rangle/\dd t=2k\langle n\rangle+k-2\gamma\langle n^2\rangle+\gamma\langle n\rangle$。
减去 $\dd\langle n\rangle^2/\dd t$ 得方差方程，稳态方差即 $\mu$。
最后，由均值方程，条件均值为
$\mathbb E[n(t)\mid n(0)]=\mu+(n(0)-\mu)e^{-\gamma t}$。两侧乘以 $n(0)-\mu$ 并在稳态下平均，就得到协方差公式。
\end{proof}
因此稳态变异系数 $\sqrt{\operatorname{Var}(n)}/\langle n\rangle=1/\sqrt\mu$，分子数越低，相对内禀涨落越大。这里平均值服从确定性线性方程是精确结果；一般非线性反应有 $\langle F(n)\rangle\ne F(\langle n\rangle)$，不能直接照搬。

\begin{figure}[htbp]\centering\includegraphics[width=.74\linewidth]{figures/birth_death.pdf}\caption{基本出生--死亡模型的稳态泊松分布。}\end{figure}
启动子随机切换 $D_{\rm off}\rightleftharpoons D_{\rm on}$ 会产生表达爆发；慢启动子切换本身就可能导致双峰分布。因此双峰直方图不构成确定性双稳态的充分证据。识别机制需结合时间序列、驻留时间、扰动响应与遗传谱系等信息。


\paragraph{为什么慢切换可以造成两个概率峰}
取二态启动子，关闭时产生率为 $k_0$、开启时为 $k_1$，两种状态共享降解率 $\gamma$，切换率为 $k_{\rm on},k_{\rm off}$。若两种切换均远慢于 $\gamma$，蛋白质常有足够时间接近当前状态的条件平衡。因此领先阶近似为
\[
P_n\simeq(1-q)\operatorname{Pois}_{k_0/\gamma}(n)
+q\operatorname{Pois}_{k_1/\gamma}(n),\qquad
q=\frac{k_{\rm on}}{k_{\rm on}+k_{\rm off}}.
\]
$q$ 来自二态链的稳态流量平衡。若两个均值相距足够远且两权重均显著，这个混合分布可以双峰，尽管每个固定启动子状态下的浓度模型都只有唯一稳定定态。该混合式是慢切换近似，有限切换速率下不是精确稳态公式。


\section{连接真实 GRN：三个必要的桥梁}
\subsection{组合调控：符号拓扑不等于启动子逻辑}
两个转录因子 $X,Y$ 结合启动子 $D$ 的四个状态为 $D,DX,DY,DXY$。在特定准平衡模型中，统计权重为
\begin{equation}1,\quad x/K_x,\quad y/K_y,\quad\omega xy/(K_xK_y),\end{equation}
$\omega$ 表示协同性。令 $Z=1+x/K_x+y/K_y+\omega xy/(K_xK_y)$，则
\begin{equation}F(x,y)=\frac{r_0+r_xx/K_x+r_yy/K_y+r_{xy}\omega xy/(K_xK_y)}{Z}.\end{equation}
同样两条激活边，由 $r_s$ 的不同可形成 AND-like、OR-like 等不同调控逻辑。仅有网络连边不能确定动力学方程。


\paragraph{四状态加权及 AND/OR 极限的推导}
写 $a=x/K_x,b=y/K_y$。准平衡时四种状态的概率是各自统计权重除以总权重 $Z$，所以 $F=\sum_s r_s p_s$，即前述有理式。$\omega=1$ 时 $Z=(1+a)(1+b)$，记 $p_x=a/(1+a),p_y=b/(1+b)$。
若只有双结合态转录，$r_{xy}=r$、其余为零，得到
$F_{\rm AND}=rp_xp_y$；若除空态以外均以速率 $r$ 转录，得到
$F_{\rm OR}=r(p_x+p_y-p_xp_y)$。
二者对 $x,y$ 均非减，却有不同的低输入响应。极限 $a,b\ll1$ 或 $\gg1$ 将连续调控近似为布尔真值表。若 $r_{xy}$ 小于单结合态速率，增加另一个因子甚至可能降低输出；因此指定某些 $r_s$ 后还要实际检查偏导数，不能只凭结合就称为激活。
\begin{figure}[htbp]\centering
\begin{tikzpicture}[node distance=20mm]
\node[gene] (d) {$D$};\node[gene,right=of d] (dx) {$DX$};
\node[gene,below=of d] (dy) {$DY$};\node[gene,right=of dy] (dxy) {$DXY$};
\draw[<->] (d)--node[above]{$X$}(dx);\draw[<->] (d)--node[left]{$Y$}(dy);
\draw[<->] (dx)--node[right]{$Y$}(dxy);\draw[<->] (dy)--node[below]{$X$}(dxy);
\node[align=left,right=15mm of dx,anchor=west] {状态权重：$1,a,b,\omega ab$\\状态转录率：$r_0,r_x,r_y,r_{xy}$};
\end{tikzpicture}
\caption{两种调控因子的启动子占据状态。图示可逆结合；协同性必须与相容的结合常数共同指定，最终输出由状态概率与各状态转录率共同决定。}
\end{figure}


\subsection{网络模体和局部功能}
真实网络中常统计自调控、前馈回路、反馈环等小拓扑是否相对某一合理随机网络零模型富集。富集不等于因果功能，也不能自动认定局部回路独立于网络其余部分。较妥当的问题是：哪些局部动力学性质在给定输入变化和网络背景下依然可观察？

\subsection{网络嵌入与有效模块}
将已分析的局部变量记为 $\bm x$，其余网络记为 $\bm z$：
\begin{equation}\dot{\bm x}=\bm f(\bm x,\bm z),\qquad\dot{\bm z}=\bm g(\bm x,\bm z).\end{equation}
线性化矩阵按块写为 $J=\begin{psmallmatrix}A&B\\C&D\end{psmallmatrix}$。即便 $A,D$ 分别稳定，耦合后的系统仍可能失稳；两块的特征值通常不足以推断整体特征值。只有在外部输入近似固定、存在明确时间尺度分离或耦合可控时，局部回路才可能构成有效的自治模块。这解释了为何最小回路有用，但不能无条件拼装成真实系统。


\paragraph{耦合失稳的反例与有效雅可比}
取 $A=D=-1$、$B=C=2$。孤立模块均为指数稳定，整体矩阵
$\begin{psmallmatrix}-1&2\\2&-1\end{psmallmatrix}$ 的特征值却是 $1,-3$，所以耦合足以使系统失稳。这是线性扰动系统的反例，不把可取负值的扰动直接解释为绝对浓度。
对于 $\lambda I-D$ 可逆的 $\lambda$，块消元给出
\[
\det(\lambda I-J)=\det(\lambda I-D)
\det\left[\lambda I-A-B(\lambda I-D)^{-1}C\right].
\]
证明是左乘 $\begin{psmallmatrix}I&B(\lambda I-D)^{-1}\\0&I\end{psmallmatrix}$ 消去右上块，得到块下三角矩阵后取行列式；消元矩阵行列式为 $1$。该表达式显示反馈项依赖频率 $\lambda$。
若外部变量的稳定弛豫显著更快，零频近似 $0=C\delta\bm x+D\delta\bm z$ 给出 $\delta\bm z=-D^{-1}C\delta\bm x$，从而 $A_{\rm eff}=A-BD^{-1}C$。即使能作时间尺度消去，所得模块也通常不是原来的 $A$。


\section{生物学案例：乳糖操纵子及其限度}
大肠杆菌 lac 系统包括 LacI 阻遏蛋白、lacZYA 转录单位以及乳糖/异乳糖相关代谢与摄取过程。lacI 具有自身的转录调控，不属于 lacZYA 同一多顺反子转录单位。异乳糖降低 LacI 的阻遏活性；CAP--cAMP 及葡萄糖相关效应调节转录和摄取。

在最简单的结合占据近似下，活性 LacI 浓度为 $R_a$ 时，解除阻遏的分数可写作 $K_d/(K_d+R_a)$，从而给出有效转录率 $F(R_a)$。进一步引入 LacY 蛋白 $y$ 与细胞内诱导物 $a$，
\begin{equation}\dot y=F(a)-\gamma_y y,\qquad\dot a=v_{\rm in}(y,L_{\rm ext})-v_{\rm out/met}(a,y).\end{equation}
这体现了 ``LacY 增强摄取 $\to$ 诱导物增加 $\to$ 阻遏解除 $\to$ LacY 增加'' 的有效正反馈，但真正的乳糖代谢、透过酶动力学、葡萄糖作用和生长稀释会改变反馈强度。特别应区分代谢性乳糖与非代谢型诱导物的实验情形，不能笼统声称 lac 网络必然双稳态。


\begin{figure}[htbp]\centering
\begin{tikzpicture}
\node[gene] (y) at (0,0) {LacY};
\node[gene] (a) at (3,0) {胞内诱导物};
\node[gene] (r) at (6,0) {活性 LacI};
\node[gene] (f) at (9,0) {lacZYA 表达};
\draw[act] (y)--node[above]{摄取}(a);
\draw[rep] (a)--(r);\draw[rep] (r)--(f);
\draw[act] (f.south) -- ++(0,-.7) -| (y.south);
\end{tikzpicture}
\caption{lac 系统的简化有效反馈关系。两次抑制形成解除阻遏，闭环总体为正；该图省略了代谢、葡萄糖作用和稀释，不据此断言双稳态。}
\end{figure}

\paragraph{从结合关系到回路符号}
单个位点近似下 $F(R_a)=r_{\rm leak}+(r_{\rm max}-r_{\rm leak})K_d/(K_d+R_a)$，当 $r_{\rm max}>r_{\rm leak}$ 时 $\partial F/\partial R_a<0$。
若有效活性阻遏蛋白 $R_a(a)$ 随诱导物下降，则链式法则给出 $\dd F/\dd a=(\partial F/\partial R_a)(\dd R_a/\dd a)>0$。
再结合 $\partial v_{\rm in}/\partial y>0$，即可识别摄取支路的正反馈。完整第二个方程的交叉导数却是
$\partial\dot a/\partial y=\partial_yv_{\rm in}-\partial_yv_{\rm out/met}$，其符号未必正；这正是代谢性诱导物与仅摄取的理想模型之间需要检查的环节。

\section{一张机制地图与自检清单}
\begin{center}\begin{tabular}{p{2.6cm}p{5.2cm}p{3.8cm}}\toprule
模型 & 可确立的典型结论 & 不应推出的结论\\\midrule
线性产生--降解 & 唯一指数稳定定态；出生--死亡离散模型为泊松稳态 & 所有表达分布均为泊松\\
一维单调自抑制 & 唯一稳定定态；无自治周期轨道 & 负反馈系统普遍不振荡\\
一维自激活 & 足够非线性时可能双稳态 & 正反馈必定双稳态\\
双组分互相调控 & 正/负回路符号影响特征值和稳定性 & 仅凭符号能决定所有行为\\
三节点负反馈 & 可发生复特征值过虚轴、支持振荡机制 & 所有负反馈环都会自振\\
前馈回路 & 可实现脉冲、滤波和特定条件下适应 & 前馈拓扑自动保证精确适应\\
\bottomrule\end{tabular}\end{center}

\noindent 阅读和复现每个模型时，依次检查：\textbf{反应规则}是什么？\textbf{消去变量}需要什么时间尺度？\textbf{定态}是否完整？\textbf{局部稳定性}与\textbf{全局动力学}是否区分？\textbf{网络结构}是否真的唯一指定调控函数？\textbf{生物功能}是否得到独立实验支持？

\section*{参考文献与延伸阅读}\addcontentsline{toc}{section}{参考文献与延伸阅读}
\begin{enumerate}[leftmargin=*]
\item Bintu, L. et al. (2005). Transcriptional regulation by the numbers: models. \emph{Current Opinion in Genetics \& Development}, 15, 116--124. DOI: \href{https://doi.org/10.1016/j.gde.2005.02.007}{\nolinkurl{10.1016/j.gde.2005.02.007}}.
\item Gardner, T. S., Cantor, C. R., \& Collins, J. J. (2000). Construction of a genetic toggle switch in Escherichia coli. \emph{Nature}, 403, 339--342. DOI: \href{https://doi.org/10.1038/35002131}{\nolinkurl{10.1038/35002131}}.
\item Elowitz, M. B. \& Leibler, S. (2000). A synthetic oscillatory network of transcriptional regulators. \emph{Nature}, 403, 335--338. DOI: \href{https://doi.org/10.1038/35002125}{\nolinkurl{10.1038/35002125}}.
\item Shen-Orr, S. S. et al. (2002). Network motifs in the transcriptional regulation network of Escherichia coli. \emph{Nature Genetics}, 31, 64--68. DOI: \href{https://doi.org/10.1038/ng881}{\nolinkurl{10.1038/ng881}}.
\item Alon, U. (2007). Network motifs: theory and experimental approaches. \emph{Nature Reviews Genetics}, 8, 450--461. DOI: \href{https://doi.org/10.1038/nrg2102}{\nolinkurl{10.1038/nrg2102}}.
\item Ma, W., Trusina, A., El-Samad, H., Lim, W. A., \& Tang, C. (2009). Defining network topologies that can achieve biochemical adaptation. \emph{Cell}, 138, 760--773. DOI: \href{https://doi.org/10.1016/j.cell.2009.06.013}{\nolinkurl{10.1016/j.cell.2009.06.013}}.
\end{enumerate}
\end{document}
